\documentclass[11pt,a4paper]{article}
\usepackage[T1]{fontenc}
\usepackage[utf8]{inputenc}
\usepackage{lmodern}
\usepackage[margin=19mm,top=17mm,bottom=18mm]{geometry}
\usepackage{amsmath,amssymb,bm}
\usepackage{microtype}
\usepackage{hyperref}
\usepackage{titlesec}
\hypersetup{hidelinks,pdftitle={Learning local structure and motion from simulation},pdfauthor={}}
\titleformat{\section}{\large\bfseries}{}{0pt}{}
\titlespacing*{\section}{0pt}{10pt}{5pt}
\titlespacing*{\paragraph}{0pt}{7pt}{0.6em}
\setlength{\parindent}{0pt}
\setlength{\parskip}{4pt}
\setlength{\abovedisplayskip}{6pt}
\setlength{\belowdisplayskip}{6pt}
\setlength{\emergencystretch}{1em}
\clubpenalty=10000
\widowpenalty=10000
\newcommand{\R}{\mathbb R}
\newcommand{\E}{\mathbb E}
\newcommand{\norm}[1]{\left\lVert #1\right\rVert}
\newcommand{\x}{\mathbf x}
\newcommand{\X}{\mathbf X}
\newcommand{\z}{\mathbf z}
\newcommand{\Sim}{\operatorname{Sim}}
\newcommand{\avgl}[1]{\left\langle #1\right\rangle_{\mathrm{long}}}
\newcommand{\avgs}[1]{\left\langle #1\right\rangle_{\mathrm{short}}}

\begin{document}
\begin{center}
{\LARGE\bfseries Learning local structure and motion from simulation}\\[4pt]
{\small Current implemented formulation}
\end{center}
\vspace{-4pt}

We learn a representation of a central atom and its local neighbourhood from
present positions and velocities. Training combines three sources of information:
structural measurements of present configurations, recorded MD futures, and
short simulated continuations with derivatives with respect to their initial
conditions. The interatomic potential is fixed. Nucleus-birth labels are used
only to evaluate the resulting representation.

\section{Starting snapshot, local input, and physical time}

\paragraph{Indices and complete starting state.}
As in the baseline, $i$ denotes a starting snapshot and $j$ the tracked central
atom. Snapshot $i$ contains $n_i$ atoms. At observation time $t$, define
\[
\begin{aligned}
\X_i(t)&=(\mathbf r_{i1}(t),\ldots,\mathbf r_{in_i}(t)),\\
\mathbf V_i(t)&=(\mathbf v_{i1}(t),\ldots,\mathbf v_{in_i}(t)),\\
\mathcal S_i(t)&=(\X_i(t),\mathbf V_i(t),\mathbf H_i(t)),
\end{aligned}
\]
where $\mathbf H_i(t)$ specifies the periodic cell. This is the physical state
used to initialize the short continuations. Thermostat settings belong to the
simulation protocol. Unlike the position-only baseline, the current method
conditions on the saved velocities rather than resampling them at the start.

\paragraph{Local observation.}
Let $\mathcal N_{ij}(t)$ contain the centre and its 79 nearest neighbours.
For $\ell\in\mathcal N_{ij}(t)$, let $\mathbf u_{ij\ell}(t)$ be the minimum-image
displacement from $j$ to $\ell$, and define the relative velocity
\[
\boldsymbol\nu_{ij\ell}(t)=\mathbf v_{i\ell}(t)-\mathbf v_{ij}(t).
\]
The geometric and phase-space observations are
\[
\x_{ij}(t)=\{\mathbf u_{ij\ell}(t)\}_{\ell\in\mathcal N_{ij}(t)},
\qquad
\boldsymbol\chi_{ij}(t)=\bigl(\x_{ij}(t),\boldsymbol\nu_{ij}(t)\bigr).
\]
The structural encoder uses positions only, with support restricted to $8$\,\AA{}
around the centre. Velocities enter through the motion coordinate defined below.
No temperature, simulation age, absolute time, species identity, or relaxed
configuration is supplied to the representation or its predictive readouts.

\paragraph{Time notation.}
The present is the recorded snapshot at physical time $t$. A horizon $\tau>0$
refers to the physical endpoint $t+\tau$. Recorded observation times are denoted
$s_k$; a future target is used only when its required time $s_k=t+\tau$ is available.
The MD integration step $\delta t_{\mathrm{MD}}$ and the saving interval $\Delta t$
are different quantities. For a trajectory saved every $m$ steps,
\[
\Delta t=m\,\delta t_{\mathrm{MD}}.
\]
The motion observations use an exact $\Delta t=0.1$\,ps grid; a source with a finer
native grid contributes its actual matching frames. No intermediate physical
frames are inferred by interpolation. The short simulator continuations
use $\delta t_{\mathrm{short}}=1$\,fs. Their step need not equal the integration
step of the archived MD trajectories.

For the learned dynamics, $\rho\geq0$ denotes elapsed physical time after the
observation, so the corresponding absolute time is $t+\rho$. Its numerical
integration step is $h=0.05$\,ps. The scales
\[
\delta t_{\mathrm{MD}},\quad \delta t_{\mathrm{short}},\quad
\Delta t,\quad h,\quad \tau
\]
therefore have distinct meanings. Neither the observation time $t$ nor a time
index is an input to a neural layer; the horizon specifies how long to evolve
the learned state.

\section{Structural initialization and a fixed structural teacher}

\paragraph{Structural representation.}
A geometry-only MACE encoder maps the local positions to
\[
\z_{ij}(t)=E_\theta(\x_{ij}(t))\in\R^d,\qquad d=128.
\]
The current encoder uses correlation order $3$ and maximum angular order $3$.
It is first trained on a broad collection of unrelaxed neighbourhoods, including
Al, Mg, Ti and Ta, with fixed material-dependent length normalization applied
before encoding. Material identity itself is not an encoder input.

\paragraph{Spatial objective.}
Two slightly perturbed views of the same neighbourhood are mapped to similar
embeddings. VICReg variance and covariance terms discourage collapse and redundant
coordinates. A structural head also predicts fixed measurements
\[
\boldsymbol\psi(\x)=
\bigl(\boldsymbol\psi_{\mathrm{BO}}(\x),
      \boldsymbol\psi_{\mathrm{TDA}}(\x)\bigr)\in\R^{298},
\]
consisting of 40 bond-order and 258 topological descriptors. The two descriptor
families have equal total weight, using normalization fitted on training data.
The objective is
\[
L_{\mathrm{spatial}}
=25L_{\mathrm{invariance}}+25L_{\mathrm{variance}}
 +L_{\mathrm{covariance}}+L_{\mathrm{descriptors}}+L_{\mathrm{masked}}.
\]
In the first initialization, $L_{\mathrm{masked}}=0$. In the second, a separate
predictor uses a partially masked neighbourhood to predict the complete-view
embedding of an exponential-moving-average teacher. The centre and its 12
nearest neighbours remain visible; other atoms are removed with probability
$0.2$. This teacher is used during spatial pretraining only.

\paragraph{A fixed reference for motion learning.}
After 12 spatial epochs, the resulting encoder is denoted $E_0$. Motion training
starts from its weights, but a separate copy of $E_0$ remains fixed and defines
all future structural targets. Thus the target representation does not change
when the trainable encoder changes. The broad spatial collection is not replayed
during motion training.

\section{A representation of present structure and motion}

\paragraph{Directional motion coordinate.}
Let $t_*=0.1$\,ps be a fixed scale that gives the motion coordinate the same units
as the structural embedding. Define
\begin{equation}
\boxed{
\mathbf h_{ij}(t)=
\begin{pmatrix}\z_{ij}(t)\\ \mathbf w_{ij}(t)\end{pmatrix},
\qquad
\mathbf w_{ij}(t)=t_*\,
D E_\theta(\x_{ij}(t))[\boldsymbol\nu_{ij}(t)]\in\R^{128}.
}
\label{eq:motion-input}
\end{equation}
Here $D E(\x)[\mathbf v]$ is a directional derivative with respect to the local
coordinates. It measures how the structural embedding responds to an
instantaneous ballistic displacement in the observed velocity direction.
It is not a separate velocity encoder or a full reconstruction of the physical
state. In particular, it does not encode the surrounding full cell or barostat
state.

\paragraph{Current numerical approximation.}
In the joint-encoder variant, the current training evaluates this coordinate as
\begin{equation}
\widetilde{\mathbf w}_{ij}(t)=
\frac{E_\theta(\x_{ij}(t)+\varepsilon_v t_*\boldsymbol\nu_{ij}(t))
      -E_\theta(\x_{ij}(t)-\varepsilon_v t_*\boldsymbol\nu_{ij}(t))}
     {2\varepsilon_v},
\qquad \varepsilon_v=0.01.
\label{eq:motion-secant}
\end{equation}
The dimensionless $\varepsilon_v$ is a differentiation amplitude.
The shifted point clouds in \eqref{eq:motion-secant} are algebraic probes of the
encoder, not saved configurations at $t\pm\varepsilon_v t_*$. Neighbour identities
are held fixed for these probes.

\paragraph{Joint and adapter variants.}
Joint training updates $E_\theta$. The cheaper adapter variant keeps $E_0$ fixed
and learns a residual map $G_\eta:\R^{128}\to\R^{128}$, initialized as the identity:
\[
\z=G_\eta(\z_0),\qquad
\mathbf w=D G_\eta(\z_0)[\widetilde{\mathbf w}_0],
\qquad \z_0=E_0(\x).
\]
The base motion coordinate $\widetilde{\mathbf w}_0$ uses
\eqref{eq:motion-secant} with $E_0$; the adapter derivative is evaluated
analytically. The adapter can reorganize information in the fixed representation,
while joint training can also change the underlying geometric features.

\section{Learned stochastic evolution and future measurements}

\paragraph{Autonomous latent dynamics.}
Starting from $\mathbf h_{ij}(t)$, evolve a stochastic latent state
\[
\mathbf H_{ij}(t,\rho)=
\bigl(\mathbf Z_{ij}(t,\rho),\mathbf W_{ij}(t,\rho)\bigr),
\qquad \mathbf H_{ij}(t,0)=\mathbf h_{ij}(t).
\]
Its continuous-time formulation is
\begin{equation}
\begin{aligned}
d\mathbf Z&=\frac{\mathbf W}{t_*}\,d\rho,\\
d\mathbf W&=\frac{\mathbf a_\vartheta(\mathbf Z,\mathbf W)}{t_*}\,d\rho
 +\frac{\operatorname{diag}\mathbf b_\vartheta(\mathbf Z,\mathbf W)}{\sqrt{t_*}}
 \,d\boldsymbol{\mathcal B}_\rho.
\end{aligned}
\label{eq:sde}
\end{equation}
Here $\boldsymbol{\mathcal B}_\rho$ is Brownian motion in elapsed time, the learned
acceleration includes damping, and the diffusion amplitudes are positive and
bounded. The coefficients depend on the latent state, not explicitly on time.
The current implementation numerically evolves this model with step $h$; its
continuous-time formulation does not imply exact dynamics between saved frames.

We use $p=1,\ldots,P$ for independent draws of the learned latent process,
separate from the physical-continuation index $b$. Training uses $P=4$ and
future evaluation uses $P=8$.

\paragraph{Two measurements of a recorded future.}
At each horizon in
\[
\mathcal T=\{0.1,\,0.5,\,3,\,12\}\ \mathrm{ps},
\]
the neighbourhood is reselected around the same atom $j$ at the actual saved
time $t+\tau$. It supplies
\begin{equation}
\mathbf Y_{ij}(t,\tau)=\phi(\x_{ij}(t+\tau))\in\R^{24},
\qquad
\mathbf A_{ij}(t,\tau)=E_0(\x_{ij}(t+\tau))\in\R^{128}.
\label{eq:long-targets}
\end{equation}
The fixed $\phi$ contains 12 smooth radial-density measurements and 12 angular
power measurements: spherical-harmonic orders $1$--$6$ in two radial shells.
Its coordinates are standardized using training observations, giving
$\widetilde{\mathbf Y}_{ij}(t,\tau)$. The second target retains the broader
structural information learned during pretraining.

These long-horizon examples contain one recorded realization for each starting
observation, rather than an average over newly generated continuations from that
same state. The learned distribution is fitted across many local observations.
Latent particles are model samples, not additional physical MD trajectories.

\paragraph{A predictive distribution.}
A decoder maps each latent endpoint to a mean and diagonal variance for the 24
standardized measurements. Writing these outputs as
$\boldsymbol\mu_{ij}^{(p)}(t,\tau)$ and
$\boldsymbol\sigma_{ij}^{2,(p)}(t,\tau)$ gives
\begin{equation}
\boxed{
\widehat p(\mathbf y\mid\boldsymbol\chi_{ij}(t);\tau)
=\frac1P\sum_{p=1}^{P}
\mathcal N\!\left(\mathbf y;
\boldsymbol\mu_{ij}^{(p)}(t,\tau),
\operatorname{diag}\boldsymbol\sigma_{ij}^{2,(p)}(t,\tau)\right).
}
\label{eq:mixture}
\end{equation}
The mixture predicts uncertainty in the chosen measurements conditional on the
local input. It does not specify a distribution over complete atomic trajectories.

Let $\avgl{\cdot}$ denote averaging over the long-future training observations.
The likelihood and structural-anchor losses are
\begin{align}
L_{\mathrm{future}}
&=\avgl{\frac1{|\mathcal T|}\sum_{\tau\in\mathcal T}
-\frac1{24}\log\widehat p\!\left(
\widetilde{\mathbf Y}_{ij}(t,\tau)\mid\boldsymbol\chi_{ij}(t);\tau\right)},
\label{eq:future-loss}\\
L_{\mathrm{anchor}}
&=\avgl{\frac1{|\mathcal T|}\sum_{\tau\in\mathcal T}
\frac1{2d}\norm{\overline{\mathbf Z}_{ij}(t,\tau)-\mathbf A_{ij}(t,\tau)}^2},
\quad
\overline{\mathbf Z}_{ij}(t,\tau)=\frac1P\sum_p\mathbf Z_{ij}^{(p)}(t,\tau).
\label{eq:anchor-loss}
\end{align}
The structural mean is therefore supervised as well as the decoded distribution.

\section{Short physical continuations and response supervision}

\paragraph{The physical experiment.}
From $\mathcal S_i(t)$, run $B=2$ short full-cell continuations to
$t+\tau_s$, where $\tau_s=0.1$\,ps:
\[
\Gamma_{ib}(t:t+\tau_s)
=\Sim_{\mathrm{short}}\!\left(\mathcal S_i(t),\xi_{ib};t,t+\tau_s\right).
\]
Both start with the same saved positions and velocities. Their independent
$\xi_{ib}$ specify thermostat noise. The simulator uses the fixed Al MEAM
potential and fixed-box Langevin dynamics integrated by BAOAB, with the source
thermostat temperature. The rest of the full cell evolves normally; the local
80-atom input is not simulated as an isolated cluster.

For the short targets, the original neighbour identities $\mathcal N_{ij}(t)$
are tracked throughout the continuation. Let
$\x_{ibj}^{\mathrm{fix}}(t,\tau_s)$ denote their centre-relative positions at
physical time $t+\tau_s$. This differs from the reselected neighbourhood used
for the recorded long-future targets.

\paragraph{Perturbations of the initial state.}
Use $r\in\{Q,V\}$ to label a position or velocity perturbation, and let $a$ be
a dimensionless amplitude. For atoms in the initial local neighbourhood,
\[
\begin{aligned}
\mathbf r_{i\ell}^{(j,r,a)}(t)
 &=\mathbf r_{i\ell}(t)+a\,\mathbf d^{Q,r}_{ij\ell},\\
\mathbf v_{i\ell}^{(j,r,a)}(t)
 &=\mathbf v_{i\ell}(t)+a\,\mathbf d^{V,r}_{ij\ell}.
\end{aligned}
\]
Both directions are zero outside the local neighbourhood. For $r=Q$, only
$\mathbf d^{Q,Q}$ is nonzero, with root-mean-square vector magnitude
$0.01$\,\AA. For $r=V$, only $\mathbf d^{V,V}$ is nonzero, with magnitude
$0.1$\,\AA/ps. Each physical direction has zero mean over its 80 atoms.
These define the full perturbed state $\mathcal S_i^{(j,r,a)}(t)$ and its
corresponding local input $\boldsymbol\chi_{ij}^{(r,a)}(t)$.

\paragraph{Fixed-encoder value and derivative targets.}
For a perturbed continuation, define the structural measurement
\[
\mathbf Y^{\mathrm{short},(r,a)}_{ibj}(t,\tau_s)
=E_0\!\left(\x_{ibj}^{\mathrm{fix}}(t,\tau_s;r,a)\right).
\]
The population value and response are
\[
\begin{aligned}
\boldsymbol\mu^{\mathrm{short}}_{ij}(t,\tau_s;r,a)
 &=\E_{\xi}\!\left[
 \mathbf Y^{\mathrm{short},(r,a)}_{ibj}(t,\tau_s)\mid\mathcal S_i(t)\right],\\
\mathbf R^{(r)}_{ij}(t,\tau_s)
 &=\left.\frac{\partial}{\partial a}
 \boldsymbol\mu^{\mathrm{short}}_{ij}(t,\tau_s;r,a)\right|_{a=0}.
\end{aligned}
\]
The current physical targets estimate these quantities by
\begin{align}
\overline{\boldsymbol\mu}^{\mathrm{short}}_{ij}(t,\tau_s)
 &=\frac1B\sum_b E_0\!\left(\x_{ibj}^{\mathrm{fix}}(t,\tau_s;0)\right),
\label{eq:teacher-value}\\
\overline{\mathbf R}^{(r)}_{ij}(t,\tau_s)
 &=\frac1B\sum_b
D E_0\!\left(\x_{ibj}^{\mathrm{fix}}(t,\tau_s;0)\right)
 \left[\mathbf T^{(r)}_{ibj}(t,\tau_s)\right],
\label{eq:teacher-response}\\
\mathbf T^{(r)}_{ibj}(t,\tau_s)
 &=\left.\frac{\partial}{\partial a}
 \x_{ibj}^{\mathrm{fix}}(t,\tau_s;r,a)\right|_{a=0}.
\nonumber
\end{align}
Here ``$;0$'' denotes the unperturbed continuation, common to both directions.
The simulator propagates $\mathbf T^{(r)}$ by differentiating its integration
steps and forces, holding each realization of thermostat noise fixed.
Thus the physical teacher uses trajectory derivatives, not differences between
the sparsely saved $0.1$\,ps frames. The interatomic potential and $E_0$ are not
optimized by the motion loss. These physical values and derivatives are generated
once and reused throughout neural training.

\newpage
\paragraph{The corresponding learned response.}
The predicted short measurement is the mean latent structural endpoint,
\[
\widehat{\boldsymbol\mu}^{\mathrm{short}}_{ij}(t,\tau_s;r,a)
=\frac1P\sum_p
\mathbf Z^{(p)}\!\left(t,\tau_s;
\boldsymbol\chi_{ij}^{(r,a)}(t)\right).
\]
The amplitude changes the encoder input; it is not an extra argument supplied
only to the decoder. The derivative must therefore describe the response of the
whole learned input-to-future map. Its current numerical estimate is
\begin{equation}
\widehat{\mathbf R}^{(r)}_{ij}(t,\tau_s)=
\frac{\widehat{\boldsymbol\mu}^{\mathrm{short}}_{ij}(t,\tau_s;r,+\varepsilon_r)
     -\widehat{\boldsymbol\mu}^{\mathrm{short}}_{ij}(t,\tau_s;r,-\varepsilon_r)}
     {2\varepsilon_r},
\qquad \varepsilon_r=0.5.
\label{eq:student-response}
\end{equation}
The same latent noise is used on both sides. At these amplitudes, the physical
input directions correspond to $0.005$\,\AA\ RMS position shifts or
$0.05$\,\AA/ps RMS velocity shifts. The teacher still represents the derivative
at $a=0$, while \eqref{eq:student-response} is a finite-difference approximation
on the neural side. Neither $a$ nor $\varepsilon_r$ denotes elapsed time.

\section{Combined training objective and matched comparisons}

Let $\avgs{\cdot}$ denote averaging over the short-response training observations,
and write $\|\mathbf v\|_d^2=d^{-1}\|\mathbf v\|^2$. The short losses are
\begin{align}
L_{\mathrm{short}}
 &=\avgs{\frac12\left\|
 \widehat{\boldsymbol\mu}^{\mathrm{short}}_{ij}(t,\tau_s;0)
 -\overline{\boldsymbol\mu}^{\mathrm{short}}_{ij}(t,\tau_s)
 \right\|_d^2},\\
L_{\mathrm{response}}
 &=\avgs{\frac14\sum_{r\in\{Q,V\}}
 \frac{\left\|\widehat{\mathbf R}^{(r)}_{ij}(t,\tau_s)
                   -\overline{\mathbf R}^{(r)}_{ij}(t,\tau_s)\right\|_d^2}
 {\max\!\left(\left\|\overline{\mathbf R}^{(r)}_{ij}(t,\tau_s)\right\|_d^2,10^{-6}\right)}}.
\end{align}
The normalization gives comparable influence to different response scales.
The implemented motion objective is
\begin{equation}
\boxed{
L_{\mathrm{motion}}=L_{\mathrm{future}}
 +0.2L_{\mathrm{anchor}}+0.1L_{\mathrm{short}}
 +\lambda_R L_{\mathrm{response}},
\qquad \lambda_R\in\{0,0.1\}.
}
\label{eq:total}
\end{equation}
The value-only control uses $\lambda_R=0$; it retains velocities, latent stochastic
dynamics, long-future supervision, and short simulator endpoint supervision.
Thus it isolates the contribution of derivative targets, rather than comparing
motion training with spatial pretraining alone.

The two spatial initializations, two adaptation regimes, and two choices of
$\lambda_R$ give eight matched fits. Each receives 1,600 motion-training updates.
The data comprise a separate collection of saved position--velocity observations
and 176 short physical response records, rather than the million-neighbourhood
spatial-pretraining collection. Their original fitting and held-out roles are
preserved. Training samples are balanced hierarchically
over original trajectory ancestors, descendant sources, and local observations;
the averages above denote this training sampling distribution.

The model is selected by future-descriptor likelihood on an internal subset of
the original training ancestors. Held-out original ancestors are not used for
fitting or checkpoint selection. No birth AP, clustering score, or crystal-label
loss enters \eqref{eq:total}.

\paragraph{What the simulator contributes.}
Recorded MD supplies targets at the longer physical horizons. The short
simulator supplies both an endpoint and the sensitivity of that endpoint to
changes in the present positions and velocities. Gradients train the neural
representation and latent dynamics against these fixed physical targets;
there is no new physical simulation inside each optimization update.

The two target sources retain their actual protocols: the short targets are
fixed-box BAOAB continuations, whereas the recorded long futures are NPT
trajectories. They also use fixed versus reselected neighbour identities.
The common latent model is trained against both, but these targets are not
assumed to be samples from an identical continuation law. Physical derivative
supervision currently reaches $0.1$\,ps only; the $3$ and $12$\,ps targets contain
recorded future values, not long-trajectory derivatives.

\section{Evaluation of future prediction and nucleus birth}

\paragraph{Future prediction.}
On held-out ancestors, evaluate the likelihood of recorded future measurements,
their mean squared prediction error, and predictive-interval coverage. Compare
with keeping the present measurements unchanged and with a distribution fitted
to training futures without the current input. Separately evaluate the short
endpoint and derivative errors. These distinguish learning future values,
learning uncertainty, and learning response to perturbations.

\paragraph{Birth prediction as a frozen readout.}
Let $T_e$ be the catalogue time of birth event $e$, and let $\tau_b>0$ be a
pre-birth lead. The input time is
\[
t_e(\tau_b)=T_e-\tau_b,\qquad t_e(\tau_b)+\tau_b=T_e,
\qquad \tau_b\in\{1,3,6\}\ \mathrm{ps}.
\]
This uses the same distinction between observation time and future horizon as
the baseline. At $t_e(\tau_b)$, extract the saved positions and velocities around a
founder atom, together with matched control centres from the same source and
time. The consolidated birth catalogue supplies the labels. No additional
neighbourhood crystallinity, component-size, or distance filter is applied.
Controls exclude catalogued birth-founder identities; they are not guaranteed
to remain liquid. Availability and deduplication checks retain the same cases
at all three leads.

With the motion model frozen, compare readouts of
\[
\z_{ij}(t_e(\tau_b)),\qquad
\mathbf h_{ij}(t_e(\tau_b)),\qquad
\bigl(\widehat{\boldsymbol\mu}^{Y}_{ij}(t_e(\tau_b),\tau_b),
       \log\widehat{\boldsymbol\sigma}^{2,Y}_{ij}(t_e(\tau_b),\tau_b)\bigr).
\]
The last vector contains the mean and marginal variance of the decoded mixture,
not a new physical continuation. The latent model is evaluated at $\tau_b$ even
when that lead is not one of its training horizons in $\mathcal T$.
The readouts use logistic regression and
boosted trees, fitted to birth labels on the original training ancestors and
selected by predictive likelihood. Original validation, calibration and test
ancestors are pooled for final evaluation. AP, AUROC, log loss and Brier score
are reported with ancestor-level uncertainty estimates and equal total weight
per physical birth event.

This is a retrospective birth-versus-matched-control assay, not an estimate of
the population nucleation rate. Its purpose is to test whether the learned
motion state or predicted future distribution adds information beyond the
structural representation, and whether derivative supervision improves that
information relative to otherwise matched value-only training.

\end{document}
